\section{Four-pole rigidity and the eventual small-cut maximum}
\label{sec:four-coefficient-rigidity}

This argument uses previously audited terminal inequalities.
It does not change their proofs or constants.
For a rank-$h$ four-pole, normalize all path and forest counts by $2^h$.
Write $a_{ij}$ for its six path counts and $u_\alpha$ for its three forest-pairing counts.
Put $S=\sum a_{ij}$ and $F=\sum u_\alpha$.
Each path coordinate and $F$ belongs to $[0,1]$.

We use the following established consequences, with absolute constants.
\begin{enumerate}
\item Every path coordinate is either one or at most $15/16$
      (Theorem~\ref{thm:terminal-path-saturation-gap}).
\item Two path coordinates sharing a terminal and exceeding $0.51$ force the third coordinate to be $O(2^{-h/2})$
      (Theorem~\ref{thm:three-pole-two-high}).
\item If $a_{12}=1$, then each of $\min(a_{13},a_{23})$ and $\min(a_{14},a_{24})$ is at most $3\,2^{-h/4}$
      (Lemma~\ref{lem:one-full-tripole-small-adjacent}).
\item A cross sum at least $2.01$ gives $F=O((h+1)2^{-h/2})$
      (Corollary~\ref{cor:linear-four-pole-cross-sum-forest}).
\item If both coordinates of an opposite pair $\alpha$ exceed $0.01$, then
      $\max(a_{ij},a_{kl})+F-u_\alpha\le1.01$ for sufficiently large $h$
      (Lemma~\ref{lem:four-pole-added-edge}).
\end{enumerate}
Suppressing unused ports gives these consequences for four-poles.
A loop produced by suppression would force the relevant path densities to be at most one half, so it does not affect the high-density cases.

\begin{lemma}[A fixed gap outside two saturated profiles]
\label{lem:four-coefficient-aggregate-rigidity}
There is an absolute $h_0$ such that every loopless normalized four-pole of rank $h\ge h_0$ with
\[
 S+F\ge3.99
\]
has either a full three-edge star or a full $K_{2,2}$ among its path coordinates.
Here an edge is full when its path count equals $2^h$.
\end{lemma}

\begin{proof}
Increase $h_0$ until every two-high small coordinate and every one-full minimum in the preceding list is at most $0.01$.
Also require the forest bound in item four to be at most $0.01$.

First suppose some cross sum is at least $2.01$.
Then $F\le0.01$, so $S\ge3.98$.
Form a graph on the four ports using edges with $a_{ij}>\tau=0.51$.
It is triangle-free, since two such adjacent edges make the third coordinate at most $0.01$.

If this high graph is $K_{2,2}$ but one edge is not full, then
\[
 S\le3+15/16+2(0.01)=3.9575<3.98.
\]
If it has three edges, it is a star or a path.
Their sums are at most $3.03$ and $3+0.51+0.02=3.53$, respectively.
Two adjacent high edges give $S\le2+0.01+3(0.51)=3.54$.
At most one high edge gives $S\le1+5(0.51)=3.55$.

The remaining case consists of two disjoint high edges.
If neither is full, then
\[
 S\le2(15/16)+4(0.51)=3.915.
\]
If one is full, item three makes two distinct cross coordinates at most $0.01$.
The other two cross coordinates are at most $0.51$.
Thus $S\le2+2(0.51)+2(0.01)=3.04$.
All alternatives contradict $S\ge3.98$.
The four high edges must therefore form a full $K_{2,2}$.

Now suppose every cross sum is less than $2.01$.
The three cross sums add to $2S$, so $S<3.015$.
The hypothesis gives $S\ge2.99$ and $F>0.975$.
Every opposite pair of coordinates has sum greater than $2.99-2.01=0.98$.

If no coordinate is full, each is at most $15/16$ and consequently greater than $0.98-15/16=0.0425$.
Apply item five to all three opposite pairings.
Their sum gives $S/2+2F\le3.03$, hence $S+4F\le6.06$.
But $S\ge2.99$ and $F>0.975$ give $S+4F>6.89$, a contradiction.

There is therefore a full coordinate, say $a_{12}=1$.
Item three gives one small coordinate in each of $\{a_{13},a_{23}\}$ and $\{a_{14},a_{24}\}$.
Each opposite coordinate exceeds $0.98-0.01=0.97$, so the saturation gap makes it full.
The two small coordinates cannot be opposite, because their opposite-pair sum exceeds $0.98$.
The three full coordinates therefore form a star.
\end{proof}

\begin{lemma}[Aligned saturated cross profiles]
\label{lem:four-coefficient-aligned-rings}
Glue two saturated four-poles across their four ports, with their full $K_{2,2}$ partitions aligned.
If either pole has no through pair, the resulting cubic graph has a cyclic cut of size three.
Otherwise the result is a square and single-rung ring.
\end{lemma}

\begin{proof}
With positive through count, the saturated four-pole classification gives a chain of square and single-rung tiles.
The two aligned chains close into the ring described in Section~\ref{sec:four-coefficient-rings}.

For completeness, the zero-through alternative in the proof of Lemma~\ref{lem:saturated-four-pole-double-cap} gives a bridge inside that pole.
It separates its left terminal pair from its right terminal pair, with positive rank on each side.
One bridge shore has two external port edges in the glued graph.
Its boundary thus consists of those two edges and the bridge.
Both global shores are connected and cyclic.
This is a cyclic three-edge cut.
\end{proof}

\begin{theorem}[Eventual exact maximum in the small-cut class]
\label{thm:four-coefficient-eventual-small-cut}
There is an absolute rank threshold $R$ such that every connected loopless cubic multigraph of rank $r\ge R$ with a cyclic cut of size at most four satisfies
\[
 c(G)\le H(r),
\]
where $H$ is the Horak--Kim count in Theorem~\ref{thm:four-coefficient-ring-maximum}.
The bound is attained for every $r\ge R$.
\end{theorem}

\begin{proof}
The established theorem for cyclic cuts of size at most three has error $O(2^{r/2})$.
That error is eventually smaller than $H(r)-2^{r-1}$.
We may therefore assume no such cut exists.

Normalize a cyclic four-edge cut as in the previously audited four-cut proof.
Its shores are connected four-poles with distinct ports and positive ranks $h\le k$.
They satisfy $r=h+k+3$.
Call the poles $X,Y$, with $Y$ of rank $k$.
In particular $k\to\infty$ uniformly as $r\to\infty$.
The exact normalized crossing count is
\[
 \mathcal B(X,Y)
 =\sum_{i<j}a_{ij}(X)a_{ij}(Y)
  +\sum_\alpha u_\alpha(Y)\bigl(F(X)-u_\alpha(X)\bigr).
\]
Every coefficient from $X$ is at most one, so $\mathcal B(X,Y)\le S(Y)+F(Y)$.

First suppose all six path coordinates of $Y$ are at most $0.51$.
The existing low-state lemma gives $S(Y)+F(Y)\le3.82$.
The asymptotic cycle bound gives $c(Y)\le(9/16)2^k$ for large $k$.
If $h\ge2$, then
\[
 \frac{c(G)}{2^{h+k}}
 \le3.82+\frac9{16}2^{-h}+2^{-k}
 \le3.960625+2^{-k}<4.
\]
If $h=1$, then $X=C_4$.
Its exact forest vector instead gives the low-state bound $\mathcal B(X,Y)\le3.7$.
Therefore
\[
 \frac{c(G)}{2^{h+k}}
 \le3.7+\frac9{32}+2^{-k}<4
\]
for sufficiently large $k$.
These bounds are uniform in the smaller rank.

Now suppose a path coordinate of $Y$ exceeds $0.51$.
The high-path cycle tradeoff gives $c(Y)\le C2^{k/2}$ for an absolute $C$.
If $S(Y)+F(Y)<3.99$, then
\[
 \frac{c(G)}{2^{h+k}}
 \le3.99+C2^{-h-k/2}+2^{-k}<4
\]
for large $k$, uniformly for every $h\ge1$.
We may assume $S(Y)+F(Y)\ge3.99$ and apply the rigidity lemma.

Suppose first that $Y$ has a full star.
Its other three path coordinates are $O(2^{-k/2})$.
If $X$ misses a full coordinate on this star, the saturation gap loses at least $1/16$ in the crossing sum.
Since $F(Y)\le1$, this gives
\[
 \mathcal B(X,Y)\le4-1/16+O(2^{-k/2}).
\]
The same normalized internal-cycle bounds make $c(G)<2^{r-1}$ for large $k$.

If $X$ is full on the same star, its three leaf-pair coordinates are $O(2^{-h/2})$.
The three full path products contribute exactly three.
The forest term is at most $F(X)F(Y)\le1$.
The remaining path products contribute $O(2^{-(h+k)/2})$.
Both poles have full paths, so their internal cycle counts are $O(2^{h/2})$ and $O(2^{k/2})$.
Consequently
\[
 c(G)\le2^{r-1}+C'2^{r/2}
\]
with an absolute $C'$, including all bounded smaller ranks.
This is eventually less than $H(r)$.

Finally suppose $Y$ has a full $K_{2,2}$.
Its remaining path coordinates are $O(2^{-k/2})$ and its forest total is $O((k+1)2^{-k/2})$.
If $X$ misses any of these four full coordinates, then
\[
 \mathcal B(X,Y)\le4-1/16+O((k+1)2^{-k/2}).
\]
The internal counts again give $c(G)<2^{r-1}$ for large $k$, uniformly in $h$.
Otherwise $X$ and $Y$ are aligned saturated cross poles.
The zero-through case in Lemma~\ref{lem:four-coefficient-aligned-rings} contradicts the absence of a cyclic three-edge cut.
The other case is a square and single-rung ring, so Theorem~\ref{thm:four-coefficient-ring-maximum} gives $c(G)\le H(r)$.

All thresholds and implied constants used above are absolute.
Their maximum supplies $R$.
The Horak--Kim rings have cyclic four-edge cuts and attain $H(r)$ on both parities.
\end{proof}

\paragraph{Status.}
This refinement has passed independent line audits by the root, path-stability, and random-host agents.
Its proof treats the smaller shore without hiding its rank in a constant.
The companion checker tests 1059 exact rational endpoint profiles; their maximum path sum outside a full $K_{2,2}$ is $1583/400=3.9575$.
The eventual rank threshold is not made effective here.
